\subsection{Environment and health}
\label{sec:revi-medio-ambiente}

Air pollution is a health risk associated with conditions such as heart
disease, stroke, pneumonia, lung cancer and respiratory infections, mainly
through exposure to particulate matter (\ce{PM2.5} and \ce{PM10}), nitrogen
dioxide (\ce{NO2}), ozone (\ce{O3}) and sulfur dioxide (\ce{SO2}). Air
quality is measured by monitoring stations that record hourly
concentrations, from which metrics such as 24-hour or annual averages are
computed depending on the pollutant, and compared against Mexican official
standards (NOMs) and the recommendations of the World Health Organization
(WHO) \parencite{oms2021}.

Some Mexican limits are less stringent than those recommended by the WHO:
for \ce{PM2.5}, the Mexican annual limit is
$10\,\mu\mathrm{g}/\mathrm{m}^{3}$, whereas the WHO recommends
$5\,\mu\mathrm{g}/\mathrm{m}^{3}$ \parencite{oms2021}. Measurement is further
complicated because not every station measures every pollutant or has
enough valid data, so sufficiency criteria are required, such as having at
least 75\,\% valid data.

The Mexican official standards set concentration limits by pollutant and
exposure period (1 hour, 8 hours, 24 hours or annual average)
\parencite{nlAirenlgobmxNormatividad, NOMCalidadAireAmbiente}. The current
limit values for ozone and particulate matter, which are used as the
exceedance criterion in this work, are given in \cref{tab:nom-limites}.
